\chapter{Projets de fin de cours}
\label{ch:capstone}

\begin{quote}
\emph{<< La meilleure fa\c{c}on d'apprendre l'IA g\'en\'erative est de construire quelque chose de r\'eel. Ces cinq projets vous m\`enent du concept \`a une application fonctionnelle. >>}
\end{quote}

% --- hook ---
Les neuf chapitres précédents ont posé chaque brique séparément. Ce dernier les rassemble. Cinq projets, chacun construit sur des données réelles et qui doit \emph{fonctionner} --- pas juste tourner dans un notebook, mais être déployé quelque part qu'un lecteur extérieur peut interroger. Le but n'est pas de produire le modèle le plus astucieux. C'est de livrer quelque chose qu'un recruteur, un collaborateur ou un domain expert peut ouvrir, comprendre, exécuter, et tenir pour fiable en quinze minutes. Si vous n'avez le temps d'en faire qu'un, choisissez celui qui répond à une question que vous, vous vous posez vraiment.
\section{Projet 1 : Application RAG sur un corpus de PDF}

\textbf{Objectif :} Construire un syst\`eme de question-r\'eponse sur une collection de documents PDF.

\textbf{Chapitres utilis\'es :} 1 (plongements), 4 (prompting), 6 (RAG).

\subsection{Sp\'ecification}

\begin{enumerate}
  \item Charger 5--10 documents PDF (articles de recherche, manuels, rapports).
  \item D\'ecouper les documents avec chevauchement (500--800 caract\`eres).
  \item Stocker les plongements dans ChromaDB avec m\'etadonn\'ees (nom de fichier, num\'ero de page).
  \item Construire une cha\^ine RAG LangChain avec Gemini comme LLM.
  \item Ajouter des citations de sources : chaque r\'eponse doit r\'ef\'erencer le document et la page.
  \item \'Evaluer avec 10 questions de test dont vous connaissez la bonne r\'eponse.
\end{enumerate}

\begin{minted}{python}
# Code de d\'epart : RAG avec suivi des sources
from langchain_community.document_loaders import PyPDFDirectoryLoader
from langchain_text_splitters import RecursiveCharacterTextSplitter
from langchain_community.vectorstores import Chroma
from langchain_community.embeddings import HuggingFaceEmbeddings
from langchain_google_genai import ChatGoogleGenerativeAI
from langchain_core.prompts import ChatPromptTemplate
from langchain_core.runnables import RunnablePassthrough
from langchain_core.output_parsers import StrOutputParser

# Charger tous les PDF d'un r\'epertoire
loader = PyPDFDirectoryLoader("./pdf_corpus/")
docs = loader.load()
print(f"Loaded {len(docs)} pages from PDFs")

# D\'ecouper
splitter = RecursiveCharacterTextSplitter(chunk_size=600, chunk_overlap=80)
chunks = splitter.split_documents(docs)

# Plonger et stocker
embeddings = HuggingFaceEmbeddings(model_name="all-MiniLM-L6-v2")
vectorstore = Chroma.from_documents(chunks, embeddings,
                                     persist_directory="./capstone_rag_db")
retriever = vectorstore.as_retriever(search_kwargs={"k": 5})

# Cha\^ine RAG avec sources
llm = ChatGoogleGenerativeAI(model="gemini-2.0-flash", temperature=0.2)

template = ChatPromptTemplate.from_template(
    """Answer the question based on the context below. Include source
references (document name, page number) for each claim.

Context:
{context}

Question: {question}

Answer (with sources):"""
)

def format_docs_with_sources(docs):
    formatted = []
    for d in docs:
        source = d.metadata.get("source", "unknown")
        page = d.metadata.get("page", "?")
        formatted.append(f"[{source}, p.{page}]\n{d.page_content}")
    return "\n\n".join(formatted)

chain = (
    {"context": retriever | format_docs_with_sources,
     "question": RunnablePassthrough()}
    | template | llm | StrOutputParser()
)

answer = chain.invoke("What are the main findings of the study?")
print(answer)
\end{minted}

\subsection{Livrables}

\begin{itemize}
  \item Notebook Jupyter fonctionnel avec pipeline complet.
  \item Tableau d'\'evaluation : 10 questions, r\'eponses attendues, r\'eponses g\'en\'er\'ees, note de justesse (1--5).
  \item R\'eflexion de 200 mots sur les limites du RAG rencontr\'ees.
\end{itemize}

% ============================================================
\section{Projet 2 : Fine-tuner un chatbot de domaine}

\textbf{Objectif :} Fine-tuner TinyLlama sur un dataset d'instructions personnalis\'e pour cr\'eer un assistant sp\'ecialis\'e.

\textbf{Chapitres utilis\'es :} 1 (tokenisation), 3 (g\'en\'eration), 5 (fine-tuning).

\subsection{Sp\'ecification}

\begin{enumerate}
  \item Choisir un domaine (cuisine, histoire africaine, programmation Python).
  \item Cr\'eer 200+ paires instruction-r\'eponse (m\'elange r\'edig\'e \`a la main et curat\'e).
  \item Fine-tuner TinyLlama avec QLoRA ($r=16$, 3 \'epoques).
  \item Comparer mod\`ele de base vs mod\`ele fine-tun\'e sur 20 questions de test.
  \item Mesurer la perplexit\'e avant et apr\`es fine-tuning.
\end{enumerate}

\begin{minted}{python}
# Aide \`a la cr\'eation de dataset
import json

def create_instruction_dataset(topic, num_examples=50):
    """Use an LLM to help generate instruction-response pairs."""
    from groq import Groq
    client = Groq()

    pairs = []
    prompt = f"""Generate {num_examples} diverse instruction-response pairs
about {topic}. Format as JSON array with keys: instruction, response.
Each response should be 2-4 sentences. Be accurate and educational."""

    response = client.chat.completions.create(
        model="llama-3.1-8b-instant",
        messages=[{"role": "user", "content": prompt}],
        temperature=0.8, max_tokens=4096,
    )

    try:
        pairs = json.loads(response.choices[0].message.content)
    except json.JSONDecodeError:
        print("Failed to parse JSON. Manual cleanup needed.")

    return pairs

# G\'en\'erer et sauvegarder
pairs = create_instruction_dataset("West African cuisine", 50)
with open("cooking_dataset.json", "w") as f:
    json.dump(pairs, f, indent=2)
print(f"Generated {len(pairs)} pairs")
\end{minted}

% ============================================================
\section{Projet 3 : Pipeline de g\'en\'eration d'images}

\textbf{Objectif :} Construire une application de g\'en\'eration d'images avec contr\^ole de style et variations.

\textbf{Chapitres utilis\'es :} 7 (mod\`eles de diffusion).

\subsection{Sp\'ecification}

\begin{enumerate}
  \item Cr\'eer un pipeline qui g\'en\`ere des images \`a partir de prompts.
  \item Impl\'ementer des pr\'esets de style (photor\'ealiste, aquarelle, art num\'erique, croquis).
  \item Ajouter image-vers-image pour des variations d'images existantes.
  \item Utiliser ControlNet pour la g\'en\'eration guid\'ee par contours.
  \item Produire une galerie de 20 images \`a travers 5 styles diff\'erents.
\end{enumerate}

\begin{minted}{python}
from diffusers import StableDiffusionPipeline, StableDiffusionImg2ImgPipeline
import torch

class ImageGenerator:
    STYLE_PROMPTS = {
        "photorealistic": "photorealistic, 8k, highly detailed, sharp focus",
        "watercolor": "watercolor painting, soft colors, artistic",
        "digital_art": "digital art, vibrant colors, detailed illustration",
        "sketch": "pencil sketch, black and white, detailed drawing",
        "oil_painting": "oil painting, rich textures, masterpiece",
    }

    def __init__(self, model_id="stabilityai/stable-diffusion-2-1-base"):
        self.pipe = StableDiffusionPipeline.from_pretrained(
            model_id, torch_dtype=torch.float16
        ).to("cuda")

    def generate(self, prompt, style="photorealistic", seed=42, **kwargs):
        styled_prompt = f"{prompt}, {self.STYLE_PROMPTS[style]}"
        generator = torch.Generator("cuda").manual_seed(seed)
        image = self.pipe(
            styled_prompt,
            negative_prompt="blurry, low quality, deformed",
            generator=generator,
            num_inference_steps=30,
            guidance_scale=7.5,
            **kwargs,
        ).images[0]
        return image

gen = ImageGenerator()
for style in ["photorealistic", "watercolor", "digital_art", "sketch"]:
    img = gen.generate("A bustling market in Cotonou", style=style)
    img.save(f"market_{style}.png")
    print(f"Generated: market_{style}.png")
\end{minted}

% ============================================================
\section{Projet 4 : Assistant de recherche multi-agents}

\textbf{Objectif :} Construire un syst\`eme multi-agents o\`u les agents collaborent pour rechercher, r\'ediger et r\'eviser un rapport.

\textbf{Chapitres utilis\'es :} 4 (prompting), 6 (RAG), 9 (agents).

\subsection{Sp\'ecification}

\begin{enumerate}
  \item \textbf{Agent chercheur} : recherche de l'information et compile des notes.
  \item \textbf{Agent r\'edacteur} : r\'edige un rapport structur\'e \`a partir des notes.
  \item \textbf{Agent r\'eviseur} : v\'erifie la justesse, la coh\'erence et l'exhaustivit\'e.
  \item Utiliser LangGraph pour orchestrer le workflow avec boucles conditionnelles.
  \item Le r\'eviseur peut renvoyer le brouillon au r\'edacteur pour r\'evision (max 2 rondes).
\end{enumerate}

\begin{minted}{python}
from langgraph.graph import StateGraph, END
from langchain_google_genai import ChatGoogleGenerativeAI
from typing import TypedDict

class ResearchState(TypedDict):
    topic: str
    research_notes: str
    draft: str
    review: str
    revision_count: int
    final_report: str

llm = ChatGoogleGenerativeAI(model="gemini-2.0-flash", temperature=0.3)

def researcher(state):
    result = llm.invoke(
        f"You are a researcher. Compile detailed notes on: {state['topic']}. "
        f"Include key facts, statistics, and recent developments."
    )
    return {"research_notes": result.content}

def writer(state):
    result = llm.invoke(
        f"You are a technical writer. Write a structured report based on:\n"
        f"{state['research_notes']}\n"
        f"{'Previous review: ' + state.get('review', '') if state.get('review') else ''}"
    )
    return {"draft": result.content,
            "revision_count": state.get("revision_count", 0) + 1}

def reviewer(state):
    result = llm.invoke(
        f"You are a critical reviewer. Review this report for accuracy, "
        f"completeness, and clarity:\n{state['draft']}\n"
        f"Respond with APPROVED if the report is good, or list specific improvements."
    )
    return {"review": result.content}

def should_revise(state):
    if state.get("revision_count", 0) >= 3:
        return "finalize"
    if "APPROVED" in state.get("review", ""):
        return "finalize"
    return "revise"

def finalize(state):
    return {"final_report": state["draft"]}

# Construire le graphe
workflow = StateGraph(ResearchState)
workflow.add_node("research", researcher)
workflow.add_node("write", writer)
workflow.add_node("review", reviewer)
workflow.add_node("finalize", finalize)

workflow.set_entry_point("research")
workflow.add_edge("research", "write")
workflow.add_edge("write", "review")
workflow.add_conditional_edges("review", should_revise,
                                {"revise": "write", "finalize": "finalize"})
workflow.add_edge("finalize", END)

app = workflow.compile()
result = app.invoke({"topic": "The impact of AI on education in Africa",
                      "revision_count": 0})
print(result["final_report"])
\end{minted}

% ============================================================
\section{Projet 5 : Benchmark d'\'evaluation}

\textbf{Objectif :} Construire un benchmark d'\'evaluation complet pour comparer des LLM.

\textbf{Chapitres utilis\'es :} 3 (g\'en\'eration), 4 (prompting), 8 (\'evaluation).

\subsection{Sp\'ecification}

\begin{enumerate}
  \item Concevoir 50 cas de test r\'epartis en 5 cat\'egories : QA factuel, raisonnement, g\'en\'eration de code, r\'esum\'e, et s\^uret\'e.
  \item \'Evaluer 3 mod\`eles (GPT-2, TinyLlama, Gemini Flash) sur tous les cas.
  \item Calculer les m\'etriques automatiques : BLEU, ROUGE, perplexit\'e.
  \item Mener une \'evaluation humaine sur un sous-ensemble de 20 cas.
  \item Produire un rapport comparatif avec tableaux et visualisations.
\end{enumerate}

\begin{minted}{python}
import json
from rouge_score import rouge_scorer
import numpy as np

class LLMBenchmark:
    def __init__(self):
        self.scorer = rouge_scorer.RougeScorer(
            ["rouge1", "rouge2", "rougeL"], use_stemmer=True
        )
        self.results = []

    def add_test_case(self, category, question, reference_answer):
        self.results.append({
            "category": category,
            "question": question,
            "reference": reference_answer,
            "model_outputs": {},
        })

    def evaluate_model(self, model_name, generate_fn):
        for case in self.results:
            output = generate_fn(case["question"])
            scores = self.scorer.score(case["reference"], output)
            case["model_outputs"][model_name] = {
                "output": output,
                "rouge1_f1": scores["rouge1"].fmeasure,
                "rouge2_f1": scores["rouge2"].fmeasure,
                "rougeL_f1": scores["rougeL"].fmeasure,
            }

    def summary(self):
        for model in list(self.results[0]["model_outputs"].keys()):
            r1 = np.mean([c["model_outputs"][model]["rouge1_f1"] for c in self.results])
            r2 = np.mean([c["model_outputs"][model]["rouge2_f1"] for c in self.results])
            rl = np.mean([c["model_outputs"][model]["rougeL_f1"] for c in self.results])
            print(f"{model}: ROUGE-1={r1:.3f}, ROUGE-2={r2:.3f}, ROUGE-L={rl:.3f}")

# Utilisation
bench = LLMBenchmark()
bench.add_test_case("factual", "What is the capital of Benin?", "The capital of Benin is Porto-Novo.")
bench.add_test_case("reasoning", "If A > B and B > C, is A > C?", "Yes, by transitivity, A > C.")
# ... ajouter d'autres cas
\end{minted}

% ============================================================
\section{Grille de notation des projets}

\begin{longtable}{p{4cm}p{2cm}p{6.5cm}}
\toprule
\textbf{Crit\`ere} & \textbf{Points} & \textbf{Description} \\
\midrule
Code fonctionnel & 30 & Le code tourne sans erreur sur Colab \\
Profondeur technique & 25 & Usage correct des techniques du cours \\
\'Evaluation & 20 & \'Evaluation quantitative et/ou qualitative \\
Documentation & 15 & Notebook clair avec explications \\
R\'eflexion & 10 & Discussion des limites et am\'eliorations \\
\midrule
\textbf{Total} & \textbf{100} & \\
\bottomrule
\end{longtable}

\begin{exercise}
\begin{enumerate}
  \item Choisissez l'un des cinq projets et impl\'ementez-le compl\`etement.
  \item R\'edigez un rapport de projet de 500 mots couvrant : objectif, approche, r\'esultats, limites et travaux futurs.
  \item Pr\'esentez votre projet en une d\'emonstration de 10 minutes montrant l'application fonctionnelle.
  \item \'Evaluez par les pairs le projet d'un autre \'etudiant en utilisant la grille ci-dessus.
  \item Proposez une sixi\`eme id\'ee de projet de fin de cours qui combine au moins 3 chapitres. R\'edigez la sp\'ecification.
\end{enumerate}
\end{exercise}

\begin{ethicsbox}
Votre projet de fin de cours sera probablement le premier syst\`eme d'IA g\'en\'erative que vous construisez. Avant de le partager :
\begin{itemize}
  \item Testez pour les sorties biais\'ees ou nocives.
  \item Ajoutez des avertissements appropri\'es (<< Contenu g\'en\'er\'e par IA >>).
  \item Ne fine-tunez pas sur des donn\'ees prot\'eg\'ees ou priv\'ees sans autorisation.
  \item Consid\'erez l'impact social : qui en b\'en\'eficie et qui pourrait \^etre l\'es\'e ?
\end{itemize}
Chaque praticien de l'IA est responsable des syst\`emes qu'il construit.
\end{ethicsbox}

% ============================================================
\section{R\'esum\'e du chapitre}

\begin{itemize}
  \item Cinq projets de fin de cours couvrent RAG, fine-tuning, g\'en\'eration d'images, syst\`emes multi-agents et \'evaluation.
  \item Chaque projet int\`egre des notions de plusieurs chapitres.
  \item Les projets sont \'evalu\'es sur code fonctionnel, profondeur technique, \'evaluation, documentation et r\'eflexion.
  \item Construisez de mani\`ere responsable : tester pour le pr\'ejudice, documenter les limites, ajouter les avertissements appropri\'es.
\end{itemize}
